\SourcePage{750}{753}
(the first term under the radical in (148.10) may be neglected in comparison
with the second); consequently, in this angular region $q$ is independent of
the amount of energy transferred. When $\vartheta\ll1$, $q$ may be either
large or small compared with $1/a_0$ (where $a_0$ is a quantity of the order
of atomic dimensions). Under the same assumption concerning the amount of
energy transferred, we have
\begin{equation}
 qa_0\sim1\quad\text{when}\quad\vartheta\sim v_0/v.
 \tag{148.13}
\end{equation}

We now return to the general formula (148.9) and consider the case of small
$q$ ($qa_0\ll1$, i.e. $\vartheta\ll v_0/v$).

In this case the exponential factors may be expanded in powers of $\mathbf q$:
$e^{-i\mathbf q\mathbf r_a}\approx1-i\mathbf q\mathbf r_a=1-iqx_a$ (the
$x$ axis is along the vector $\mathbf q$). On substituting this expansion in
(148.9), the terms containing 1 vanish by virtue of the orthogonality of the
wave functions $\psi_0$ and $\psi_n$, and we obtain
\begin{equation}
 d\sigma_n=8\pi\left(\frac{e}{\hbar v}\right)^2\frac{dq}{q}|\langle n|d_x|0\rangle|^2
 =\left(\frac{2e}{\hbar v}\right)^2|\langle n|d_x|0\rangle|^2\frac{d\Omega}{\vartheta^2},
 \tag{148.14}
\end{equation}
where $d_x=e\sum_a x_a$ is the component of the atomic dipole moment. We see
that the scattering cross-section (at small $q$) is determined by the squared
modulus of the dipole-moment matrix element for the transition corresponding
to the change of the atomic state\footnote{Usually the quantity of physical
interest is the cross-section $d\sigma_n$, summed over all orientations of the
atomic moment in the final state and averaged over its orientations in the
initial state. After this summation and averaging, the square
$|\langle n|d_x|0\rangle|^2$ no longer depends on the direction of the $x$
axis.}.

It may happen, however, that the dipole-moment matrix element for the given
transition vanishes identically because of the selection rules (a forbidden
transition). The expansion of $\exp(-i\mathbf q\mathbf r_a)$ must then be
continued to the next term, and we obtain
\begin{equation}
 d\sigma_n=2\pi\left(\frac{e^2}{\hbar v}\right)^2
 \left|\left\langle n\left|\sum_a x_a^2\right|0\right\rangle\right|^2q\,dq.
 \tag{148.15}
\end{equation}

We now consider the opposite limiting case of large $q$ ($qa_0\gg1$). Large
$q$ means that the momentum transferred to the atom is large compared with
the initial intrinsic momenta of the atomic electrons. It is physically clear
in advance that in this case the atomic electrons may be treated as free, and
the collision with the atom as an elastic collision of the incident electron
with atomic electrons initially at rest. This also follows from the general
formula (148.9). At large $q$, the integrand in the matrix element contains
rapidly oscillating factors $\exp(-i\mathbf q\mathbf r_a)$, and the integral
can fail to be small only if $\psi_n$ contains the same factor.

\SourcePage{751}{754}
Such a function $\psi_n$ corresponds to an ionized atom from which an electron
has been ejected with momentum $-\hbar\mathbf q=\mathbf p-\mathbf p'$, fixed
simply by momentum conservation, just as in a collision between two free
electrons.

In a collision involving a large momentum transfer, the two electrons (the
incident and atomic electrons) may both emerge with speeds of comparable
magnitude. The exchange effects associated with the identity of the colliding
particles, which were not taken into account in the general formula (148.9),
then become important. The scattering cross-section for fast electrons with
exchange included is given by (137.9); that formula applies in the frame in
which one of the electrons was at rest before the collision. For fast
electrons, the cosine in the last term of (137.9) may be replaced by unity.
Multiplying also by the number $Z$ of electrons in the atom, we obtain the
cross-section for collision of an electron with an atom in the form
\begin{equation}
 d\sigma=4Z\left(\frac{e^2}{mv^2}\right)^2
 \left(\frac1{\sin^4\vartheta}+\frac1{\cos^4\vartheta}
 -\frac1{\sin^2\vartheta\cos^2\vartheta}\right)\cos\vartheta\,d\Omega.
 \tag{148.16}
\end{equation}

In this formula it is convenient to express the scattering angle in terms of
the energy acquired by the electrons after the collision. As is known, when a
particle of energy $E=mv^2/2$ collides with a particle of the same mass at
rest, the energies of the particles after the collision are
$\varepsilon=E\sin^2\vartheta$ and $E-\varepsilon=E\cos^2\vartheta$. To obtain
the cross-section referred to the interval $d\varepsilon$, we express
$d\Omega$ in terms of $d\varepsilon$ according to
$\cos\vartheta\,d\Omega=2\pi\cos\vartheta\sin\vartheta\,d\vartheta
=(\pi/E)d\varepsilon$. Substitution in (148.16) gives the final formula
\begin{equation}
 d\sigma_\varepsilon=\pi Ze^4\left[\frac1{\varepsilon^2}+\frac1{(E-\varepsilon)^2}
 -\frac1{\varepsilon(E-\varepsilon)}\right]\frac{d\varepsilon}{E}.
 \tag{148.17}
\end{equation}
If one of the energies $\varepsilon$ or $E-\varepsilon$ is small compared with
the other, only one of the three terms in this formula (the first or the
second) is important. This corresponds to the fact that, when the two electron
energies differ greatly, the exchange effect is unimportant and we must return
to the usual Rutherford formula\footnote{For the collision of a positron with
an atom, the exchange effect is altogether absent, and the Rutherford formula
$d\sigma_\varepsilon=(\pi Ze^4/E)d\varepsilon/\varepsilon^2$ applies for all
$q\gg1/a_0$.}.

Integration of the differential cross-section over all angles (or,
equivalently, over $dq$) gives the total cross-section $\sigma_n$ for a
collision that excites the given atomic state. The dependence of $\sigma_n$
on the incident electron's speed is essentially connected with whether the
atomic dipole-moment matrix element for the corresponding transition is
present or absent. Suppose first that

\SourcePage{752}{755}
this element is nonzero. At small $q$, the cross-section $d\sigma_n$ is then
given by (148.14), and we see that the integral over $dq$ diverges
logarithmically as $q$ decreases. In the region of large $q$, by contrast, the
cross-section (for a specified energy transfer $E_n-E_0$) decreases
exponentially as $q$ increases, because of the already noted rapidly
oscillating factor in the integrand of the matrix element in (148.9). Thus the
region of small $q$ gives the main contribution to the integral over $dq$, and
we may restrict the integration to the range from the minimum value
$q_{\min}$ in (148.11) to some value $\sim1/a_0$. The result is
\begin{equation}
 \sigma_n=8\pi\left(\frac{e}{\hbar v}\right)^2|\langle n|d_x|0\rangle|^2
 \ln\left(\beta_n\frac{v\hbar}{e^2}\right),
 \tag{148.18}
\end{equation}
where $\beta_n$ is a dimensionless constant that cannot be calculated in
general\footnote{We assume that $E_n-E_0$ is of the order of the atomic-electron
energy $\varepsilon_0$. For large energy transfers
($E_n-E_0\sim E\gg\varepsilon_0$), formulae (148.14), (148.18) are inapplicable
in any event, since the dipole-moment matrix element becomes very small and
one cannot retain only the first term of the expansion in $q$.}.

If, on the other hand, the dipole-moment matrix element vanishes for the given
transition, the integral over $dq$ converges rapidly both at small $q$ (as is
clear from (148.15)) and at large $q$. In this case the principal region for
the integral is $q\sim1/a_0$. No general quantitative formula can be obtained
here, and we conclude that $\sigma_n$ is inversely proportional to the square
of the speed:
\begin{equation}
 \sigma_n=\frac{\mathrm{const}}{v^2}.
 \tag{148.19}
\end{equation}
This follows directly from the general formula (148.9), according to which
$d\sigma_n$ at $q\sim1/a_0$ is proportional to $v^{-2}$.

Let $d\sigma_r$ be the cross-section for inelastic scattering into a given
solid-angle element, irrespective of the state to which the atom passes. To
find it, expression (148.9) must be summed over all $n\ne0$, that is, over all
atomic states (of both the discrete and continuous spectra) except the ground
state. We exclude from consideration both large and extremely small angles
and assume that $(v_0/v)^2\ll\vartheta\ll1$. Then, according to (148.12), $q$
is independent of the energy transferred\footnote{The sum in (148.9) also
extends over states with $E_n-E_0\gg\varepsilon_0$, for which (148.12) does not
hold. Transitions involving a large energy transfer have, however, a
comparatively small cross-section, and these terms contribute little to the
sum. The condition $\vartheta\ll1$ permits exchange effects to be neglected.}.

\SourcePage{753}{756}
The last circumstance makes it easy to calculate the total cross-section for
inelastic collisions, that is, the sum
\begin{equation}
\begin{split}
 d\sigma_r&=\sum_{n\ne0}d\sigma_n
 =8\pi\left(\frac{e^2}{\hbar v}\right)^2\sum_{n\ne0}
 \left|\left\langle n\left|\sum_a e^{-i\mathbf q\mathbf r_a}\right|0\right\rangle\right|^2\frac{dq}{q^3}\\
 &=\left(\frac{2e^2}{mv^2}\right)^2\sum_{n\ne0}
 \left|\left\langle n\left|\sum_a e^{-i\mathbf q\mathbf r_a}\right|0\right\rangle\right|^2\frac{d\Omega}{\vartheta^4}.
\end{split}
 \tag{148.20}
\end{equation}

To do this, note that for any quantity $f$, the rule for matrix multiplication
gives
\[
 \sum_n|f_{0n}|^2=\sum_nf_{0n}(f_{0n})^*=\sum_nf_{0n}(f^+)_{n0}=(ff^+)_{00}.
\]
The summation here is over all $n$, including $n=0$. Therefore
\begin{equation}
 \sum_{n\ne0}|f_{0n}|^2=\sum_n|f_{0n}|^2-|f_{00}|^2=(ff^+)_{00}-|f_{00}|^2.
 \tag{148.21}
\end{equation}

Applying this relation to $f=\sum_a e^{-i\mathbf q\mathbf r_a}$, we obtain
\begin{equation}
 d\sigma_r=\left(\frac{2e^2}{mv^2}\right)^2
 \left\{\left\langle\left|\sum_a e^{-i\mathbf q\mathbf r_a}\right|^2\right\rangle
 -\left|\left\langle\sum_a e^{-i\mathbf q\mathbf r_a}\right\rangle\right|^2\right\}
 \frac{d\Omega}{\vartheta^4},
 \tag{148.22}
\end{equation}
where $\langle\ldots\rangle$ denotes averaging over the ground state of the
atom (that is, taking the diagonal matrix element 00). By definition, the mean
value $\left\langle\sum e^{-i\mathbf q\mathbf r_a}\right\rangle$ is the atomic
factor $F(q)$ of the atom in its ground state. In the first term in braces we
may write
\[
 \left|\sum_{a=1}^Z e^{-i\mathbf q\mathbf r_a}\right|^2
 =Z+\sum_{a\ne b}e^{i\mathbf q(\mathbf r_a-\mathbf r_b)}.
\]
Thus we find the general formula
\begin{equation}
 d\sigma_r=\left(\frac{2e^2}{mv^2}\right)^2
 \left\{Z-F^2(q)+\left\langle\sum_{a\ne b}e^{i\mathbf q(\mathbf r_a-\mathbf r_b)}\right\rangle\right\}
 \frac{d\Omega}{\vartheta^4}.
 \tag{148.23}
\end{equation}

This formula is greatly simplified at small $q$, when it may be expanded in
powers of $q$ ($v_0/v\ll qa_0\ll1$, corresponding to angles
$(v_0/v)^2\ll\vartheta\ll v_0/v$). Rather than carry out the expansion in
formula (148.23), it is more convenient to start anew
